% !PRINTED 280
\weberpage{280}
For $\varrho=\gamma^{12}$ we may choose an arbitrary primitive root of 11, e.g.\ 8. Then $\varrho^{-1}=7$ and
\[
\varrho-\varrho^{-1}=1,\quad \varrho^2-\varrho^{-2}=4,
\]
hence
\weq{13}{
\Theta = \begin{pmatrix} 3, & 0 \\ 0, & 4 \end{pmatrix}, \
\psi = \begin{pmatrix} 0, & 1 \\ -1, & 0 \end{pmatrix}, \
\chi = \begin{pmatrix} -3, & 1 \\ 1, & 3 \end{pmatrix}.
}
With this it is thus established which kinds of divisors can occur in the group $L_p$, and it has also been shown that all these divisors are actually present. The question of how one obtains, for each type, all the divisors that are possible at all, we have left aside here. We wish to make only the following remarks about it.

If one transforms a group $G$ that has been found by any substitution of the group $L_p$, one obtains a conjugate divisor. But if one transforms by an imaginary substitution of the group $E_p$, it can happen that the transformed group of $G$ nevertheless becomes real and is not conjugate to $G$ within $L_p$ (although both are conjugate divisors of $E_p$). Thus from the octahedral and the icosahedral group one obtains in general a second, non-conjugate one, if one transforms
with $\begin{pmatrix} \varepsilon, & 0 \\ 0, & \varepsilon^{-1} \end{pmatrix}$. For the tetrahedral group this yields a new group only when $p\equiv\pm 1\ (\mathrm{mod}\ 8)$. The second group results from the formula:
\weq{14}{
\begin{pmatrix} \varepsilon, & 0 \\ 0, & \varepsilon^{-1} \end{pmatrix}
\begin{pmatrix} a, & b \\ c, & d \end{pmatrix}
\begin{pmatrix} \varepsilon^{-1}, & 0 \\ 0, & \varepsilon \end{pmatrix}
=
\begin{pmatrix} a, & bN \\ cN^{-1}, & d \end{pmatrix}
\quad (N=\varepsilon^2).
}
In order to exhibit these circumstances in the examples $p=5,7,11$, it is expedient to transform the group (6) found above for $p=7$ in such a way that the substitution $\chi$, which is of the third order, receives the normal form $S$. One finds this transformation easily, and obtains for this case:
\[
\begin{pmatrix} 1, & 1 \\ -2, & -1 \end{pmatrix}
\begin{pmatrix} 0, & 2 \\ 3, & 1 \end{pmatrix}
\begin{pmatrix} -1, & -1 \\ 2, & 1 \end{pmatrix}
=
\begin{pmatrix} 3, & 0 \\ 0, & 5 \end{pmatrix}
= \chi'
\]
\weq{15}{
\begin{pmatrix} 1, & 1 \\ -2, & -1 \end{pmatrix}
\begin{pmatrix} 2, & 1 \\ 3, & 2 \end{pmatrix}
\begin{pmatrix} -1, & -1 \\ 2, & 1 \end{pmatrix}
=
\begin{pmatrix} 1, & -2 \\ -1, & 3 \end{pmatrix}
= \Theta'
}
\[
\begin{pmatrix} 1, & 1 \\ -2, & -1 \end{pmatrix}
\begin{pmatrix} 2, & 2 \\ 1, & -2 \end{pmatrix}
\begin{pmatrix} -1, & -1 \\ 2, & 1 \end{pmatrix}
=
\begin{pmatrix} -3, & -3 \\ 1, & 3 \end{pmatrix}
= \psi',
\]
%CONT>>
